\SourceLeaf{155}
welche kennen zu lernen erst für den ein lebhaftes Interesse haben kann, der
die Methoden schon zu gebrauchen versteht. Ich pflege drei Begründungsarten
vorzutragen:

\begin{enumerate}
\item[1)] eine bloss auf Principien der Zweckmässigkeit basirte, die sich, sehr
einleuchtend, leicht machen lässt,
\item[2)] die in der Th. M. C. C. gelehrte Anknüpfungsart an die
Wahrscheinlichkeitsrechnung,
\item[3)] die davon \textls{durchaus verschiedene} Anknüpfungsart an die
Wahrscheinlichkeitsrechnung, welche in der Theoria Combinationis
Observationum vorgetragen, und nach meiner Überzeugung die ausschliesslich
einzige zulässige ist.
\end{enumerate}

Für jeden, der mit dieser Lehre noch ganz unbekannt ist, halte ich diese
Reihefolge für die zweckmässigste. Fremde Darstellungen kenne ich aber
keine, die ich empfehlen könnte. Was aber No. 1) betrifft, so kann Herr
\textsc{Warensdorf} diess aus unzähligen Büchern lernen (was das A. B. C. der Sache
betrifft) und es ist ziemlich einerlei, welches er wählt, wenn er nur das,
was fast bei allen schlecht ist, nachher Gelegenheit nimmt bei sich zu
verbessern. So z. B. legen diese Bücher keinen hinlänglichen Druck darauf,
dass man auch dann, wenn die beobachteten Grössen schon linearische
Functionen von den unbekannten Elementen sind (wie bei der Pendellänge
$x+y\sin\varphi^2$), man doch nicht diese Elemente als die unbekannten
Grössen der Aufgabe betrachten, sondern dazu die an schon bekannte
\textls{möglichst genäherte} Werthe anzubringenden Correctionen wählen soll, u. s. w.

. . . . . . . . . . . .
